\chapter{Yannakakis: \emph{tree-walk}}
\label{cap:yannakakis-tree}

\begin{intuicion}
Yannakakis chain (capítulo anterior) sirve para hipótesis cuya
\emph{shape} es una cadena lineal. Para hipótesis \emph{shape
graph} ---estrellas, árboles, cualquier hipergrafo
\(\alpha\)-acíclico no-cadena--- el join tree ya no es una línea
sino un árbol propio. El \emph{tree-walk} enumerator camina ese
árbol manteniendo un \emph{binding} (mapa de variables a vértices)
y, en cada nodo del árbol, decide qué \emph{bag} consultar según
qué variable comparte con su padre. Este capítulo explica el
algoritmo, muestra un ejemplo con una estrella, y discute el
contraste con el chain enumerator.
\end{intuicion}

\section{El setup: variables, bindings, join tree}

Para una hipótesis \emph{shape graph}, los pasos no se relacionan
posicionalmente sino vía \emph{variables} explícitas:

\begin{lstlisting}
User c -[Auth]-> Host h1; c -[Auth]-> Host h2; c -[Auth]-> Host h3
\end{lstlisting}

\noindent
Aquí \(c\) es una \texttt{VarId} compartida entre los tres pasos,
mientras \(h_1, h_2, h_3\) son \texttt{VarId}s independientes. La
hipótesis tiene 4 variables y 3 pasos.

El \emph{binding} es un \texttt{Vec<Option<StrId>>} indexado por
\texttt{VarId}; arranca \texttt{None} en todas las posiciones, y
el tree-walk lo va llenando.

\section{El join tree para \emph{shape graph}}

El planner (capítulo~\ref{cap:planner}) construye el join tree con
GYO + BFS desde una hiperarista \emph{root}. Para la estrella de
arriba:

\begin{figure}[H]
\centering
\begin{tikzpicture}[
    every node/.style={draw=accent, thick, rectangle, rounded corners=1mm,
                       minimum size=8mm, fill=accentlight, font=\small}
]
\node (s0) at (0, 1) {step 0};
\node (s1) at (-2, -0.5) {step 1};
\node (s2) at (2, -0.5) {step 2};

\draw[accent, thick] (s0) -- node[pos=0.5, fill=white, draw=none, font=\scriptsize]{$c$} (s1);
\draw[accent, thick] (s0) -- node[pos=0.5, fill=white, draw=none, font=\scriptsize]{$c$} (s2);
\end{tikzpicture}
\caption{Join tree para
\texttt{c -[Auth]-> h1; c -[Auth]-> h2; c -[Auth]-> h3}. Tres
nodos (uno por step), conexión etiquetada con la \emph{shared
variable} con el padre (todos comparten \(c\) con step 0, que es
la root).}
\label{fig:tree-walk-jointree}
\end{figure}

\subsection{La estructura \texttt{JoinTree}}

\begin{lstlisting}[language=Rust]
pub struct JoinTree {
    pub nodes: Vec<usize>,            // step indices, en orden de visita BFS
    pub parent: Vec<Option<usize>>,   // parent[i]: índice en `nodes` del padre
    pub shared_var: Vec<VarId>,       // variable compartida con el padre
}

pub struct StepBag {
    pub rows: Vec<StepRow>,                // todas las aristas del bag
    pub by_origin: HashMap<StrId, Vec<usize>>, // index: origin → indices en rows
    pub by_dest:   HashMap<StrId, Vec<usize>>, // index: dest   → indices en rows
}
\end{lstlisting}

\noindent
Cada \texttt{StepBag} tiene dos índices auxiliares: \texttt{by\_origin}
y \texttt{by\_dest}. Esto permite consultas eficientes en cualquier
sentido (a diferencia del chain enumerator, que sólo necesita
\texttt{by\_origin}).

\section{El algoritmo}

\begin{algorithm}[H]
\caption{Tree-walk enumerator (versión simplificada)}
\KwIn{\texttt{JoinTree}, \texttt{bags: Vec<StepBag>}, \texttt{binding} (todo \texttt{None})}
\KwOut{Vec<HuntResult>}
\(\mathrm{recurse}(\mathrm{node\_idx} = 0)\)\;
\BlankLine
\Comment{\texttt{recurse(node\_idx)}}
\If{\(\mathrm{node\_idx} = |\mathrm{nodes}|\)}{
    \texttt{emit}(binding resuelto a Vec<Arc<str>>)\;
    \Return\;
}
\(\mathrm{step} \leftarrow \mathrm{nodes[node\_idx]}\)\;
\(\mathrm{edge} \leftarrow \mathrm{hypergraph.edges[step]}\)\;
\Switch{\((\mathrm{binding[edge.origin]}, \mathrm{binding[edge.dest]})\)}{
    \Case{\((\mathrm{Some}(o), \_)\)}{candidates \(\leftarrow \mathrm{by\_origin[step][o]}\)}
    \Case{\((\mathrm{None}, \mathrm{Some}(d))\)}{candidates \(\leftarrow \mathrm{by\_dest[step][d]}\)}
    \Case{\((\mathrm{None}, \mathrm{None})\)}{candidates \(\leftarrow \mathrm{bags[step].rows}\) (root)}
}
\For{\(\mathrm{row} \in \mathrm{candidates}\)}{
    intentar bind(edge.origin, row.origin) y bind(edge.dest, row.dest)\;
    \(\mathrm{recurse}(\mathrm{node\_idx} + 1)\)\;
    \texttt{unbind} si correspondía\;
}
\end{algorithm}

\section{Por qué el switch (origin\_bound, dest\_bound)}

La diferencia clave con el chain enumerator: en una cadena sabemos
que \(B_{i+1}.\mathrm{origin} = B_i.\mathrm{dest}\) por la convención
posicional, así que el lookup es siempre \texttt{by\_origin}. En
una estrella o árbol, la \emph{shared variable} con el padre puede
aparecer como \texttt{origin} o como \texttt{dest} de la arista del
step actual.

\begin{ejemplo}[Estrella revisitada]
Hipótesis:
\texttt{c -[Auth]-> h1; c -[Auth]-> h2; c -[Auth]-> h3}.

\begin{itemize}
    \item Step 0: arista \texttt{c -[Auth]-> h1}. Su \texttt{origin}
    es \(c\), su \texttt{dest} es \(h_1\). Como root, ambas
    variables están \texttt{None}; iteramos \texttt{bags[0].rows}
    enteros. Cada iteración bind \(c \leftarrow \mathrm{row.origin}\)
    y \(h_1 \leftarrow \mathrm{row.dest}\).
    \item Step 1: arista \texttt{c -[Auth]-> h2}. Su \texttt{origin}
    es \(c\) (ya bound desde step 0), su \texttt{dest} es \(h_2\)
    (None). El switch dispatcha al caso
    \((\mathrm{Some}(o), \mathrm{None})\): candidates =
    \texttt{by\_origin[1][\(o\)]} con \(o = \mathrm{binding[c]}\).
    Iteramos esos rows, bind \(h_2 \leftarrow \mathrm{row.dest}\).
    \item Step 2: análogo a step 1, bind \(h_3\).
    \item En el caso base (\(\mathrm{node\_idx} = 3\)), emit el path
    completo (\(c, h_1, h_2, h_3\)).
\end{itemize}
\end{ejemplo}

\subsection{Caso \((\mathrm{None}, \mathrm{Some}(d))\)}

¿Cuándo aparece este caso? Cuando un step tiene su \texttt{dest}
bound pero no su \texttt{origin}. Por ejemplo, una hipótesis
``árbol invertido'':

\begin{lstlisting}
Host h -[Inbound]-> User u1; Host h2 -[Inbound]-> u1
\end{lstlisting}

\noindent
Acá \(u_1\) (en el dest) es la shared variable. Step 0 bind \(h\)
y \(u_1\); step 1 tiene \(\mathrm{dest} = u_1\) (already bound)
y \(\mathrm{origin} = h_2\) (None). Lookup vía
\texttt{by\_dest[1][\(u_1\)]}.

\subsection{Caso \((\mathrm{Some}, \mathrm{Some})\)}

Step donde tanto origin como dest ya están bound. El motor consulta
si la arista \emph{específica} existe (lookup determinista, no
iteración). Si existe, recursa; si no, fail-fast.

\section{Bindings y unbinding}

\begin{lstlisting}[language=Rust]
fn bind(
    binding: &mut Vec<Option<StrId>>,
    var: VarId,
    sid: StrId,
) -> bool {
    match binding[var.0] {
        None => { binding[var.0] = Some(sid); true }
        Some(existing) => existing == sid,  // ya bound: chequear igualdad
    }
}

fn unbind(
    binding: &mut Vec<Option<StrId>>,
    var: VarId,
    was_just_bound: bool,
) {
    if was_just_bound {
        binding[var.0] = None;
    }
}
\end{lstlisting}

\noindent
Cada \texttt{bind} en el \texttt{recurse} loop es ``si la variable
ya estaba bound, chequear que coincide; si no, bind y marcar para
unbind al volver''. El \texttt{unbind} restaura sólo lo que se
acaba de bind ---\emph{no} desbinde lo que ya estaba bound desde
arriba---.

\section{Análisis}

Igual que Yannakakis chain, la complejidad es
\(O(|\E_{\text{relevantes}}| + |Q_{\text{out}}|)\). El reducer
(no mostrado en detalle aquí) reduce los bags como en chain; el
tree-walk recorre el join tree sin visitas duplicadas.

\section{Cuándo se dispatcha}

GraphHunter dispatcha al tree-walk cuando:

\begin{enumerate}
    \item La hipótesis tiene \texttt{shape == HypothesisShape::Graph}.
    \item El planner verifica \(\alpha\)-aciclicidad
    (capítulo~\ref{cap:hipergrafos}) ---debería ser cierto si la DSL
    valida correctamente, pero el chequeo es defensivo.
    \item El feature flag \texttt{GRAPHHUNTER\_YANNAKAKIS} está
    on (default).
\end{enumerate}

\noindent
Si la hipótesis es shape graph pero \emph{cíclica} (\(K_3\),
\(K_4\)), Yannakakis no aplica; el planner cae a LFTJ-K\(_4\) (si
es \(K_4\)) o al DFS legacy.

\section{Resumen}

\begin{itemize}
    \item Tree-walk enumerator es la versión de Yannakakis para
    shape graph: walk recursivo sobre el join tree, con
    binding/unbinding.
    \item En cada step, el lookup se decide por
    \texttt{(origin\_bound, dest\_bound)}: dispatch a
    \texttt{by\_origin}, \texttt{by\_dest}, o iteración completa
    de bag (root).
    \item Complejidad lineal en input + output.
    \item Aplica para hipótesis \(\alpha\)-acíclicas no-cadena
    (estrellas, árboles).
\end{itemize}

\begin{ejercicios}
\begin{enumerate}
    \item \starred{} Para la hipótesis
    \texttt{User c -[Auth]-> Host h1; c -[Auth]-> Host h2}
    en un grafo donde \(c\) tiene 100 vecinos \texttt{Host}, ¿cuántos
    matches enumera el tree-walk? ¿Cuál es la cota teórica?

    \item Construí el join tree para la hipótesis
    \texttt{Process p1 -[Spawn]-> Process p2 -[Spawn]-> Process p3; p1 -[Spawn]-> Process p4}.
    ¿Qué shape es?

    \item \starred{} ¿Cuándo aparece el caso
    \((\mathrm{Some}, \mathrm{Some})\) en el switch? Diseñá una
    hipótesis donde aparezca y explicá qué chequea.

    \item Si en vez de mantener dos índices (\texttt{by\_origin} y
    \texttt{by\_dest}) por bag, mantuviéramos sólo
    \texttt{by\_origin}, ¿qué cambia? ¿En qué casos sería más
    lento?

    \item Compará tree-walk con DFS bien-podado para una hipótesis
    de estrella \(c \to (h_1, h_2, h_3, h_4, h_5)\). ¿Cuál gana?
    ¿Por qué?
\end{enumerate}
\end{ejercicios}
